% ============================================================================
%  第 5 章 求解算法 —— body/05-planning-algorithm.tex
%  职责：给出算法结构（含伪代码）、关键机制、复杂度与最优性。
% ============================================================================

\section{\termTD 风险感知路径规划算法}
\label{sec:planning}

% ---------------------------------------------------------------------------
\subsection{总体思路}

直接在三维节点上再展开一维时间，会把搜索图规模放大 $T$ 倍。
\methodnamelatex 的核心是把时间从\textbf{图拓扑}中剥离，作为
\textbf{标签状态}沿路径传播，从而在不膨胀搜索图的前提下支持任意时刻的
在线风险查询。

% ---------------------------------------------------------------------------
\subsection{标签化传播}

对节点 $v$，维护标签 $\ell = (\gcost, \tlabel, \mathrm{parent})$，
其中 $\tlabel$ 为到达时刻；松弛时按
\refeqn{eq:state-def} 更新时间：
\begin{equation}
  \tlabel(v') = \tlabel(v) + \Delta t(v, v'),
  \qquad
  \gcost(v') = \gcost(v) + \riskcost\bigl((v, \tlabel(v)) \to (v', \tlabel(v'))\bigr).
  \label{eq:label-prop}
\end{equation}

% ---------------------------------------------------------------------------
\subsection{累积危险率的加法形式}

为避免长路径上概率乘积导致的数值下溢，改以累积危险率传播：
\begin{equation}
  H(v') = H(v) + \int_{\tlabel(v)}^{\tlabel(v')} \hazard(u) \, \dd{u},
  \qquad
  \survival = \exp(-H).
  \label{eq:cumulative-hazard}
\end{equation}

% ---------------------------------------------------------------------------
\subsection{算法流程}

\begin{algorithm}[t]
\caption{\methodnamelatex：时间依赖风险感知最短路搜索}
\label{alg:td-risk-astar}
\begin{algorithmic}[1]
  \Input 时空图 $G = (\statespace, \edgeset)$，起点 $(v_0, t_0)$，终点 $v_d$，权重 $\{w_i\}$
  \Output 最优路径 $\pi^{*}$
  \State $\mathrm{OPEN} \gets \{(v_0, \gcost{=}0, \tlabel{=}t_0)\}$，$\mathrm{CLOSED} \gets \emptyset$
  \While {$\mathrm{OPEN} \neq \emptyset$}
    \State $\ell \gets \argmin\{\gcost + \heur(\cdot) \mid \ell \in \mathrm{OPEN}\}$
    \State 从 $\mathrm{OPEN}$ 取出 $\ell$
    \If {$\ell.v = v_d$}
      \State \textbf{return} 回溯标签链得到 $\pi^{*}$
    \EndIf
    \For {每个后继 $v' \in \mathcal{N}(\ell.v)$}
      \State $\tlabel' \gets \ell.\tlabel + \Delta t(\ell.v, v')$
      \State $\riskcost' \gets \textsc{RiskQuery}(\ell.v, v', \tlabel')$       \Comment 在线查询动态风险张量
      \State $\ell' \gets (\gcost{=}\ell.\gcost + \riskcost',\; \tlabel{=}\tlabel',\; \mathrm{parent}{=}\ell)$
      \If {$\ell'$ 被 $\mathrm{CLOSED}(v')$ 中的标签支配} \Comment 见第 \ref{sec:dominance} 节
        \State \textbf{continue}
      \EndIf
      \State 更新 $\mathrm{OPEN}(v')$ 与 $\mathrm{CLOSED}(v')$
    \EndFor
  \EndWhile
\end{algorithmic}
\end{algorithm}

% ---------------------------------------------------------------------------
\subsection{支配剪枝}
\label{sec:dominance}

\begin{definition}[标签支配]
标签 $\ell_1$ 支配 $\ell_2$，当且仅当
$\gcost_1 \le \gcost_2$ 且 $\tlabel_1 \le \tlabel_2$，且至少一处严格不等。
\end{definition}

\begin{theorem}[剪枝安全性（单标签覆盖）]
在 \cref{thm:monotone} 的单调代价结构下，被支配标签不可能导出优于保留标签的解。
\end{theorem}

% \todo{补充剪枝正确性的完整证明，并与 \cref{thm:monotone} 的前提条件对齐。}

% ---------------------------------------------------------------------------
\subsection{复杂度分析}

设平均出度为 $b$、路径长度为 $K$、支配检查代价为 $O(\log q)$（$q$ 为标签数上界），
则时间复杂度为 $O(bK \log q)$，空间复杂度为 $O(|\nodeset| \, q)$；
在 $q$ 有界的情形下与传统三维 \basestatic \ 同阶。
